\documentclass[english,hebrew]{article}
\usepackage[HE8,T1]{fontenc}
\usepackage[utf8]{inputenc}
\setcounter{secnumdepth}{-2}
\setlength{\parindent}{0bp}
\usepackage{amsmath}
\usepackage{amsthm}
\usepackage{amssymb}

\makeatletter
\@ifundefined{date}{}{\date{}}
\usepackage{graphicx}
\usepackage{geometry}
\newcommand{\limi}{\underset{n\rightarrow\infty}{\lim}}
\newcommand{\limxz}{\underset{x\rightarrow x_{0}}{\lim}}
\newcommand{\limfxz}{\underset{x\rightarrow x_{0}}{\lim}f(x)}
\newcommand{\limfxm}{\underset{x\rightarrow x_{0}^{-}}{\lim}f(x)}
\newcommand{\limfxp}{\underset{x\rightarrow x_{0}^{+}}{\lim}f(x)}
\newcommand{\limxm}{\underset{x\rightarrow x_{0}^{-}}{\lim}}
\newcommand{\limxp}{\underset{x\rightarrow x_{0}^{+}}{\lim}}
\newcommand{\sqrm}{M_{m\times m}(\mathbb{F})}
\newcommand{\seqa}{(a_{n})_{n=1}^{\infty}}
\newcommand{\seqx}{(x_{n})_{n=1}^{\infty}}
\newcommand{\funcdr}{f:D\rightarrow\mathbb{R}}
\newcommand{\der}{\limxz\frac{f(x)-f(x_{0})}{x-x_{0}}}
\usepackage[all]{pdfprivacy}

\makeatother

\usepackage{babel}
\begin{document}
\title{משפט הערך הממוצע של לגראנז' הוא מקרה פרטי של משפט טיילור עם השארית
בצורת לגראנז'}

\maketitle
משפט טיילור עם השארית בצורת לגראנז':

יהיו \L{$a,b\in\mathbb{R}$} כך ש- \L{$a<b$} ויהי \L{$n\in\mathbb{N}\cup\left\{ 0\right\} $}.

תהי \L{$f$} פונקציה הגזירה \L{$n$} פעמים בקטע \L{$\left[a,b\right]$}
כך שהנגזרת ה-\L{$n$} שלה רציפה והגזירה \L{$n+1$} פעמים בקטע \L{$\left(a,b\right)$}.

אז קיימת נקודה \L{$c\in\left(a,b\right)$} כך ש-\L{%
\[
R_{n,f,a}\left(b\right)=f(b)-P_{n,f,a}\left(b\right)=\frac{f^{\left(n+1\right)}(c)}{\left(n+1\right)!}\left(b-a\right)^{n+1}
\]
}%

וזה בפרט נכון עבור \L{$n=0$}:

אם \L{$f$} היא פונקציה רציפה בקטע \L{$\left[a,b\right]$} וגזירה
בקטע \L{$\left(a,b\right)$}, אז קיימת נקודה \L{$c\in\left(a,b\right)$}
כך ש-\L{%
\[
f(b)-P_{0,f,a}\left(b\right)=\frac{f^{\left(0+1\right)}(c)}{\left(0+1\right)!}\left(b-a\right)^{1}
\]
}%

נשים לב:\L{%
\[
\frac{f^{\left(0+1\right)}(c)}{\left(0+1\right)!}\left(b-a\right)^{1}=f'(c)\cdot\left(b-a\right)
\]
}%
\L{%
\[
P_{0,f,a}\left(b\right)=\underset{k=0}{\overset{0}{\sum}}\frac{f^{(k)}(a)}{k!}\left(b-a\right)^{k}=\frac{f(a)}{1}\left(b-a\right)^{0}=f(a)
\]
}%
 בסך הכל קיבלנו\L{%
\[
f(b)-f(a)=f'(c)\cdot\left(b-a\right)
\]
}%

כלומר, קיימת נקודה \L{$c\in\left(a,b\right)$} כך ש-\L{%
\[
f'(c)=\frac{f(b)-f(a)}{b-a}
\]
}%
 וזה בדיוק משפט הערך הממוצע של לגראנז'. 
\end{document}
